\section{Finite convolution support at a prime simple factor}
\label{app:atomic-convolution}

This appendix proves the finite local theorem used at the width-bearing prime
currents.  The theorem statement is about support in the semiring
$\N[\mathbf F_p]$.  Matching coordinates and characters occur only inside the
proof.

\subsection{Fibre counts and convolution}

For a finite map $f:B\to S$ to a finite abelian group, define
\[
\mathsf{Card}_S(f)
:=\sum_{b\in B}[f(b)]
=\sum_{s\in S}|f^{-1}(s)|[s]
\in\N[S].
\]
Then
\[
\operatorname{supp}(\mathsf{Card}_S(f))=\operatorname{im}(f).
\tag{AC1}\label{eq:card-support}
\]
If $f_i:B_i\to S$ are finite families and
$\Sigma(b_i)=\sum_i f_i(b_i)$, finite Fubini gives
\[
\mathsf{Card}_S(\Sigma)=*_{i}\mathsf{Card}_S(f_i).
\tag{AC2}\label{eq:fubini-convolution}
\]
Thus for two-point slots $B_i=\{x_i,y_i\}$ the branch-sum map is surjective
exactly when every coefficient of
\[
*_{i}([x_i]+[y_i])
\]
is positive.

\subsection{A simultaneous energetic matching}

Let $X\subset\mathbf F_p$, $|X|=N$, and for $t\ne0$ put
\[
w_t(x,y)=\left\|\frac{t(x-y)}p\right\|^2,
\]
where $\|\cdot\|$ is circular distance to the nearest integer.

\begin{lemma}[Difference-energy shell]\label{lem:difference-energy-shell}
There is an absolute $c_0>0$ such that
\[
\sum_{x\ne y\in X}w_t(x,y)
\ge c_0\frac{N^4}{p^2}
\qquad(t\ne0).
\]
\end{lemma}

\begin{proof}
For one nonzero residue $d$, the ordered equation $t(x-y)=d$ has at most $N$
solutions.  Hence one unsigned circular-distance level has multiplicity at
most $2N$.  Among the $N(N-1)$ ordered pairs the smallest possible total
square distance is therefore obtained by filling the lowest distance levels,
with at most $2N$ pairs at each level.  Summing the first $\gg N$ squares
$j^2/p^2$ gives the claimed $\gg N^4/p^2$ bound.
\end{proof}

If $N$ is odd, discard one arbitrary point and write $X'$ for the remaining
set and $N'=N-1$; if $N$ is even put $X'=X$ and $N'=N$.  Reapply
Lemma~\ref{lem:difference-energy-shell} to $X'$ itself.  On the nontrivial
large-$N$ locus, $N'\ge N/2$, so every lower bound below expressed using $N'$
is the corresponding bound using $N$ up to an absolute constant.  The slots
constructed in $X'$ are of course also slots in $X$.  For the matching
argument only, we now rename $(X',N')$ as $(X,N)$; thus $N$ is even in the
remainder of this subsection.

Let $\Omega$ be the finite set of perfect matchings of $X$.  Every unordered
pair occurs in the same number of matchings, so
Lemma~\ref{lem:difference-energy-shell} gives
\[
\overline S_t
:=\frac1{|\Omega|}\sum_{M\in\Omega}\sum_{e\in M}w_t(e)
\ge c_1\frac{N^3}{p^2}
\tag{AC3}\label{eq:matching-average}
\]
for an absolute $c_1>0$.

\begin{lemma}[Finite matching concentration]\label{lem:matching-concentration}
There is an absolute $c_2>0$ such that the fraction of $M\in\Omega$ with
\[
S_t(M):=\sum_{e\in M}w_t(e)<\overline S_t/2
\]
is at most
\[
\exp\!\left(-c_2\frac{N^5}{p^4}\right).
\]
\end{lemma}

\begin{proof}
Reveal a matching by repeatedly pairing the first remaining vertex.  At every
node let $\mu$ be the arithmetic mean of the final energy over all leaves
below it.  Two children which pair the current vertex $a$ with $y$ and $z$ are
coupled by the switching bijection
\[
(a,y),(z,u)\longleftrightarrow(a,z),(y,u).
\]
Only two edge weights change.  Since $0\le w_t\le1/4$, the child conditional
means lie in an interval of one absolute length $b$.

The parent mean is the arithmetic average of its child means.  Hence the child
deviations have mean zero and range length at most $b$.  The elementary finite
Hoeffding inequality
\[
\frac1m\sum_{j=1}^m e^{\lambda z_j}
\le e^{C\lambda^2b^2}
\quad
\left(\sum_jz_j=0,\ \max z_j-\min z_j\le b\right)
\]
follows from convexity of the exponential on a bounded interval.  Iterating
this inequality through at most $N/2$ levels yields
\[
\frac1{|\Omega|}\sum_{M\in\Omega}
 e^{\lambda(S_t(M)-\overline S_t)}
\le e^{C'N\lambda^2}.
\]
Use $\lambda<0$ and optimize at
$|\lambda|\asymp\overline S_t/N$.  Together with
\eqref{eq:matching-average} this gives
\[
\frac{\#\{M:S_t(M)<\overline S_t/2\}}{|\Omega|}
\le
\exp\!\left(-c_2\frac{\overline S_t^2}{N}\right)
\le
\exp\!\left(-c_3\frac{N^5}{p^4}\right).
\]
\end{proof}

There is therefore an absolute $A_0$ such that, whenever
\[
\frac{N^5}{p^4}\ge A_0\log p,
\tag{AC4}\label{eq:matching-threshold}
\]
finite subadditivity over the $p-1$ nonzero $t$ gives one matching $M_*$ for
which
\[
S_t(M_*)\ge c_4\frac{N^3}{p^2}
\qquad(t\ne0)
\tag{AC5}\label{eq:simultaneous-energy}
\]
simultaneously, for an absolute $c_4>0$.

\subsection{Exact-cardinality selection}

Fix $1\le M\le|M_*|$.  Partition the edges of $M_*$ into exactly $M$ nonempty
chunks, each of size at most a constant multiple of $N/M$.  Let $\mathcal S$
be the finite Cartesian product choosing one edge from each chunk.  For
$A\in\mathcal S$ put
\[
\Phi_A:=*_{\{x_e,y_e\}\in A}([x_e]+[y_e])
\in\N[\mathbf F_p].
\]
Its augmentation is $2^M$.

For the nontrivial additive character $\chi_t$,
\[
\frac{|\chi_t(\Phi_A)|}{2^M}
=
\prod_{e\in A}
\left|\frac{\chi_t(x_e)+\chi_t(y_e)}2\right|.
\]
The elementary cosine estimate gives
\[
\left|\frac{\chi_t(x)+\chi_t(y)}2\right|
\le e^{-c_5w_t(x,y)}.
\]
Since $w_t$ is bounded, averaging over one chunk and using the chord bound for
$e^{-c_5u}$ gives
\[
\frac1{|C|}\sum_{e\in C}e^{-c_5w_t(e)}
\le
\exp\!\left(-c_6\frac{\sum_{e\in C}w_t(e)}{|C|}\right).
\]
Multiplying over chunks and applying \eqref{eq:simultaneous-energy} yields
\[
\frac1{|\mathcal S|}\sum_{A\in\mathcal S}
\frac{|\chi_t(\Phi_A)|}{2^M}
\le
\exp\!\left(-c_7\frac{N^2M}{p^2}\right).
\tag{AC6}\label{eq:chunk-character}
\]

Hence there is an absolute $A_1$ such that, whenever
\[
M\ge A_1\frac{p^2}{N^2}\log p,
\tag{AC7}\label{eq:atomic-threshold}
\]
summing \eqref{eq:chunk-character} over $t\ne0$ shows that the finite
arithmetic average of
\[
\frac1{2^M}\sum_{t\ne0}|\chi_t(\Phi_A)|
\]
is $<1$.  Some exact-$M$ selection therefore satisfies
\[
\sum_{t\ne0}|\chi_t(\Phi_A)|<2^M.
\tag{AC8}\label{eq:augmentation-dominance}
\]

Character orthogonality gives for every $s\in\mathbf F_p$
\[
p\,[s]\Phi_A
=2^M+
\sum_{t\ne0}\chi_t(-s)\chi_t(\Phi_A).
\]
By \eqref{eq:augmentation-dominance} the right-hand side is strictly positive.
Thus every coefficient of $\Phi_A$ is positive, and
\eqref{eq:card-support} gives a surjective branch-sum map.

We have proved:

\begin{theorem}[Dense convolution-support cover]\label{thm:atomic-convolution}
There are absolute constants $A_0,A_1>0$ such that the following holds.  Let
$X\subset\mathbf F_p$ with $|X|=N$.  If
\[
N^5/p^4\ge A_0\log p,
\]
then every integer $M$ satisfying
\[
M\ge A_1(p^2/N^2)\log p,
\qquad M\le\lfloor N/2\rfloor
\]
admits $M$ disjoint two-point slots in $X$ whose branch-sum map onto
$\mathbf F_p$ is surjective.
\end{theorem}

The constants may be enlarged once to absorb the one-point reduction when
$N$ is odd.  The matching, chunks and characters have no consumer after this
theorem.
